%!TEX root = ../atomic_jsc.tex
% LTeX: language=en
\section{Preliminaries}
\label{sec:preliminaries}

%\subsection{Well-quasi-orders}
\arka{shift just before existence?}

\subsection{Partial orders, ordinals, well-founded sets, and well-quasi-ordered sets.}
\AP We assume basic familiarity with partial orders, well-founded sets, and ordinals.
We use the notation $\intro*\om$ for the first infinite ordinal (that is, $(\N, \leq)$), and write $X \intro*\ordplus Y$ for the lexicographic sum of two partial orders $X$ and $Y$
(i.e., $z < z'$ if either $z\in X$ and $z'\in Y$, or if $z < z'\in X$ or $z < z'\in Y$).
Similarly, the notation $X\intro*\ordtimes Y$ denotes the product of two partial orders equipped with the lexigographic ordering
(i.e., $(x_1, y_1) \leq (x_2, y_2)$ if either $x_1 < x_2$, or $x_1 = x_2$ and $y_1 \leq y_2$).
We also use the usual notations for finite ordinals, writing $\intro*\ordfin{n}$ for the finite ordinal of size $n$.
For instance, $\om \ordplus \ordfin{1}$ is the total order $\N \uplus\set{+\infty}$, where $+\infty$ is the new largest element.

\AP In order to guarantee the termination of the algorithms presented in this
paper, a key ingredient is the notion of \intro{well-quasi-ordering}
(WQO), which consists of sets $(X, \leq)$ such that every infinite sequence
$\seqof{x_i}[i \in \N]$ of elements of $X$ contains a pair $i < j$ such that
$x_i \leq x_j$. Examples of \kl{well-quasi-orderings} include finite sets with
any ordering, or $\N \times \N$ with the product ordering. We refer the reader
to \cite{SCSC17} for a comprehensive introduction to \kl{well-quasi-orderings}
and their applications in computer science.
\arka{Higman and Kruskal? Maybe in appendix?}
%
\subsection{Polynomials, monomials, divisibility.}
\arka{this subsection will be merged with equivariant grobner bases}
\AP We use the notation $\poly{\K}{\Indets}$ for the ring of polynomials with coefficients from a field $\K$ and indeterminates from a set $\Indets$,
and $\mon{\Indets}$ for the set of monomials in $\poly{\K}{\Indets}$.
Letters $p,q,r$ denote polynomials,
$\monelt,\monelt[n]$ denote monomials, and $a,b,\alpha,\beta$ denote coefficients in $\K$.
\arka{reserve one of them for variables or check instances where $a,b$ are used for variables}

A classical example of a \kl{WQO} is the set of monomials $\mon{\Indets}$,
endowed with the \kl{divisibility} relation $\intro*\divleq$ when $\Indets$
is finite. We recall that a monomial $\monelt[m]$ \intro{divides} a monomial
$\monelt[n]$ if there exists a monomial $\monelt[l]$ such that $\monelt[m]\monelt[l] = \monelt[n]$. In this case, we write $\monelt[m]
\reintro*\divleq \monelt[n]$.

\AP Let $\varleq$ be a total ordering on $\Indets$.
Using this order,
we define the \intro{co-lexicographic} (colex) ordering on monomials as follows:
$\monelt[n] \intro*\revlexlt \monelt[m]$ if there exists an
indeterminate $x \in \Indets$ such that $\monelt[n](x) < \monelt[m](x)$,
and such that for every $y \in \Indets$,
if $x \varlt y$ then $\monelt[n](y) = \monelt[m](y)$. Note that if $\monelt[n] \divleq \monelt[m]$,
then in particular $\monelt[n] \revlexleq \monelt[m]$. 

\AP We can now use the \kl{reverse lexicographic} ordering to identify particular elements in a given polynomial. Namely, for a polynomial $p \in \poly{\K}{\X}$, we define
the \intro{leading monomial} $\intro*\lm(p)$ of $p$ as the largest monomial
appearing in $p$ with respect to the \kl{revlex ordering}, and the
\intro{leading coefficient} $\intro*\lc(p)$ of $p$ as the coefficient of
$\lm(p)$ in $p$. We then define the \intro{leading term} $\intro*\lt(p)$ of
$p$ as the product of its \kl{leading monomial} and its \kl{leading
coefficient},
\arka{push to later}
and the \intro{characteristic monomial} $\intro*\cm(p)$ of $p$ as
the product of its \kl{leading monomial} and all the indeterminates appearing
in $p$.
We also define the \intro(monomial){domain} of $\monelt[m]$ as the set
$\intro*\dom(\monelt[m])$ of indeterminates $x \in \X$ such that $\monelt[m](x) \neq
0$. Because the coefficients and monomial in question are highly dependent on
the ordering $\varleq$, we allow ourselves to write $\lm[\Indets](p)$ to
highlight the precise ordered set of variables that was used to compute the
\kl{leading monomial} of $p$. We extend $\dom$ from 
monomials to polynomials by defining $\dom(p)$ as 
the union of the \kl(monomial){domains} of all monomials appearing in $p$.

\begin{remark}\label{rem:other-monomial-orderings}
In the case of a finite set of indeterminates, one can choose any total ordering
on $\mon{\Indets}$, as long as it contains the \kl{divisibility}
quasi-ordering, and is compatible with the product of monomials.\footnote{This
is often called a \emph{monomial ordering}, see \cite{CLO15}.} In our case,
having an infinite number of indeterminates, we rely on a connection between
$\lm(p)$ and $\dom(p)$: $\dom(p) \subseteq \dwset{\dom(\lm(p))}$, where
$\dwset{S}$ is the downwards closure of a set $S \subseteq \Indets$, i.e. the
set of all indeterminates $x \in \Indets$ such that $y \leq x$ for some $y \in
S$. This means that the \kl{leading monomial} encodes a \emph{global property}
of the polynomial, and it will be crucial in our termination arguments. This is
already at the core of the classical \emph{elimination theorems} \cite[Chapter 3, Theorem
2]{CLO15}.
\end{remark}

\subsection{Ideals and Gr\"{o}bner Bases}
\AP An \intro{ideal} $\idl$ of
$\poly{\K}{\X}$ is a non-empty subset of $\poly{\K}{\X}$ that is closed under
addition and multiplication by elements of $\poly{\K}{\X}$. Given a set $H
\subseteq \poly{\K}{\X}$, we denote by $\intro*\IdlGen{H}$ the ideal generated
by $H$, i.e. the smallest ideal that contains $H$. The \intro{ideal membership
problem} is the following decision problem: given a polynomial $p \in
\poly{\K}{\X}$ and a set of polynomials $H \subseteq \poly{\K}{\X}$, decide
whether $p$ belongs to the ideal $\IdlGen{H}$ generated by $H$. We know that
this problem is decidable when $\X$ is finite, and that it is even
$\EXPTIME$-complete \cite{MAME82}. The classical approach to the \kl{ideal
membership problem} is to use the \kl{Gröbner basis} theory, developed
in the 1970s by Buchberger~\cite{BUCH76}. 
A set $\Basis$ of polynomials is called a \intro{Gröbner basis} of
an ideal $\idl$ if, $\IdlGen{\Basis} = \idl$ and for every polynomial $p \in
\idl$, there exists a polynomial $q \in \Basis$ such that $\lm[\Indets](q)
\divleq \lm[\Indets](p)$.

Given a \kl{Gröbner basis} $\Basis$ of an ideal $\idl$, and a polynomial $p$,
it suffices to iteratively reduce the \kl{leading monomial} of $p$ by
subtracting multiples of elements in $\Basis$, until one cannot apply any
reductions. If the result is $0$, then $p$ belongs to $\idl$, and otherwise it
does not. 

\begin{example}
  Let $\Indets \defined \set{x,y,z}$ with $z < y < x$. The
  set $\Basis \defined \set{x^2y - z, x^2 - y}$ is not 
  a \kl{Gröbner basis} of the ideal $\idl$ it generates,
  because the polynomial $p \defined y^2 - z$ belongs to $\idl$
  but its \kl{leading monomial} $y^2$ is not divisible by 
  $\lm(x^2y - z) = x^2y$ nor by $\lm(x^2 - y) = x^2$.
\end{example}

%\para{Computability assumptions.}
%\AP We say that the action is
%\intro{effectively oligomorphic} if:
%%
%\begin{enumerate}
%\item It is \intro{oligomorphic}, i.e.\ for every $n \in \N$,
%$\Indets^n$ is \kl{orbit finite},
%\item There exists an algorithm that decides whether two elements $\vec{x},
%\vec{y} \in \Indets^*$ are in the same orbit under the action of 
%$\group$ on $\Indets^*$.
%\item There exists an algorithm which, on input $n\in\N$, outputs a set
%$A \subseteqfin \Indets^n$ such that $|A\cap U| = 1$ for every orbit
%$U \subseteq \Indets^n$.
%\end{enumerate}
%In particular, $\Indets$ itself is \kl{orbit finite} under the action of $\group$.
%
%
%A group action $\group \actson \X$ is said to be \intro(ord){compatible}
%with an ordering $\leq$ on $\X$ if for all $\gelem \in \group$ and $x,y \in
%\X$, we have $x \leq y$ if and only if $\gelem \cdot x \leq \gelem \cdot y$.
%Let us point out that in this case, $\revlexleq$ is also \kl(ord){compatible} with
%the action of $\group$ on $\mon{\X}$, i.e. for all $\gelem \in \group$ and
%monomials $\monelt[m], \monelt[n] \in \mon{\X}$, we have $\monelt[m] \revlexleq
%\monelt[n]$ if and only if $\gelem \cdot \monelt[m] \revlexleq \gelem \cdot
%\monelt[n]$.
%Our \intro{computability assumptions} on the tuple $(\Indets, \group,
%\leq)$ are that $\group$ acts \kl{effectively oligomorphic} on
%$\Indets$, and that its action is \kl(ord){compatible} with the ordering $\leq$
%on $\Indets$.
%
%\begin{example}
%  \label{ex:computability-assumptions}
%  Let $\Indets \defined \Q$ and $\group$ be the group of all
%  order preserving bijections of $\Q$.
%  Then, $\group$ acts \kl{effectively oligomorphically} on $\Indets$,
%  and its action is \kl(ord){compatible} with the ordering of $\Q$ by definition. 
%\end{example}
%
%Note that under our \kl{computability assumptions}, the set of polynomials
%$\poly{\K}{\Indets}$ is also \kl{effectively oligomorphic} under the action of
%$\group$ on $\Indets$ when restricted to polynomials with bounded degree.
%This is because a polynomial $p \in \poly{\K}{\Indets}$
%can be seen as an element of $(\K \times \Indets^{\leq d})^n$ where $n$ is the
%number of monomials in $p$, and $d$ is the maximal degree of a monomial
%appearing in $p$. Note that the set of all polynomials $\poly{\K}{\Indets}$ is
%not \kl{orbit finite}, precisely because the orbit of a polynomial $p$ under 
%the action of $\group$ cannot change the degree of $p$, and because there are 
%polynomials of arbitrarily large degree.
%
\subsection{Monoid actions, equivariance, orbit-finiteness and effective oligomorphism} 
\AP
In this section we define basic notions regarding monoid actions and define the concept of effective local oligomorphism.
As we show in \arka{refer to sec},
the latter condition will allow us to execute the equivariant Buchberger's algorithm (defined in \arka{refer to sec}).


A \intro{monoid} $\monoid$ is a set equipped with a binary operation that is associative, and has an identity element.
In our setting, we are interested in polynomial rings of the form $\poly{\K}{\Indets}$ where $\Indets$ is an infinite set of indeterminates that are equipped with an injective \intro{monoid action} $\monoid \actson \Indets$.
This means that there is monoid homomorphism $\psi_{\monoid}$ from $\monoid$ to the monoid of one to one endofunctions on $\Indets$.
Obviously, the same monoid $\monoid$ can act differently on a set $\Indets$.
However, usually the action $\psi_{\monoid}$ will be clear from the choice of $\monoid$ and $\Indets$ and we will use $\monoid\actson\Indets$ to specify this choice.
For each $\gelem \in \monoid$, we use $\gelem\cdot x$ to denote $\psi_{\monoid}(\gelem)(x)$.
Usually, we will be interested in actions $\monoid\actson\Indets$ where $\Indets$ comes with some additional structure and $\monoid$ is the monoid of all structure preserving injective endofunctions/group of all structure preserving bijections on $\Indets$.
For instance, $\inc{\OrderAtoms}\actson\OrderAtoms$ denote the obvious action of the order preserving one-to-one endofunctions on the structure $\OrderAtoms$ of rational numbers with order.
In \arka{cite sec} we give some more examples of such actions.  

\arka{shift to computability assumptions? Just after writing down theorem?
Computability follows from the rule of thumb that equivariant functions can be evaluated on orbit-finite sets.
We show how this rule of thumb formalises in our setting.
Maybe even formalise it?
Sets with atoms? Takes a lot of formalisms.
We take a shortcut. Input = a fixed numbers of (Tuple,string in binary). output similar.
But which kind of sets. maybe only lower levels in the hierarchy? bounded branching? heriditary finite sets or their orbits}
The action $\monoid\actson\Indets$ extends to actions $\monoid\actson\Indets^n$ for $n\in\N$ in the obvious way:
$\gelem\cdot(x_1,\dots,x_n) = (\gelem\cdot x_1,\dots,\gelem\cdot x_n)$ for $\gelem\in\monoid$ and $x_1,\dots,x_n\in\Indets$.
A subset $S \subseteq \Indets^n$ is called \intro{$\monoid$-equivariant} if for all $\gelem \in \monoid$ and $f \in S$, we have $\gelem \cdot f\in S$.
For example, $\setof{(x,y)\in\OrderAtoms^2}{x < y \in \OrderAtoms}$ is an $\inc{\OrderAtoms}$-equivariant subset of $\OrderAtoms^2$.
However, the subset $\setof{x \in \OrderAtoms}{x \text{ is an integer}}$ is not an $\inc{\OrderAtoms}$-equivariant subset of $\OrderAtoms$.

We denote $\intro*\orbit[\monoid]{P}$ for the set of all elements $\gelem \cdot p$ for $\gelem \in \monoid$ and $p \in P$.
When $P$ is a singleton set ${p}$,
we write $\orbit[\monoid]{p}$ instead of $\orbit[\monoid]{\{p\}}$.
An \kl{equivariant set} is said to be \intro{orbit finite} if it is the
union of finitely many \intro{orbits} under the action of $\monoid$.
Equivalently, an \reintro{orbit finite set}
is a set of the form $\orbit[\group]{P}$ for some finite set $P$.
%
\begin{definition}\label{def:loc oligo}
The action $\monoid\actson\Indets$ is called \intro{locally oligomorphic} if for every $x_1,\dots,x_n\in\Indets$,
the subset $\setof{(y_1,\dots,y_n)}{y_k\in\orbit[\monoid]{x_k}}$ is \kl{orbit-finite}.
The action $\monoid\actson\Indets$ is \intro{effectively locally oligomorphic} if:
\begin{enumerate}
\item it is \kl{locally oligomorphic},
\item there exists an algorithm to check whether a tuple $(a_1,\dots,a_n)\in\Indets^n$ is inside $\orbit[\monoid]{(b_1,\dots,b_n)}$ for some $(b_1,\dots,b_n)\in\Indets^n$, and
\item there exists an algorithm that takes elements $a_1,\dots,a_n\in\Indets$ as input and outputs a finite set $P$ such that $\setof{(y_1,\dots,y_n)}{y_k\in\orbit[\monoid]{x_k}} = \orbit[\group]{P}$.
\end{enumerate}
\end{definition}
%
\arka{maybe we need slightly different where tuples are chosen as input?}
%
The following lemma implies that \kl{local oligomorphicity} is a generalisation of the well-known notion of \kl{oligomorphicity} for group actions:
a group action $\group\actson\Indets$ is called \intro{oligomorphic} if sets of the form $\Indets^n$ is orbit-finite for every $n$.
This lemma will be followed by examples and non-examples of locally oligomorphic monoid actions. 
%
\begin{lemma}\label{lem:oligo loc oligo}
Let $\group\actson\Indets$ be a \kl{group action}.
\begin{enumerate}
\item If $\group\actson\Indets$ is \kl{oligomorphic} then it is \kl{locally oligomorphic}.
\item If $\Indets$ is \kl{orbit-finite} then $\group\actson\Indets$ is \kl{oligomorphic} if and only if it is \kl{locally oligomorphic}.
\end{enumerate}
\end{lemma}
%
\begin{lemma}\label{lem:orbit part}
Let $\group\actson\Indets$ be a group action.
For $(a_1,\dots,a_n),(b_1,\dots,b_n)\in\Indets^*$,
\[
(a_1,\dots,a_n)\in\orbit[\group]{(b_1,\dots,b_n)}
\quad\iff\quad
(b_1,\dots,b_n)\in\orbit[\group]{(a_1,\dots,a_n)}
\]
\end{lemma}
%
\arka{if this lemma is not used anywhere else, make it a claim and push it inside the proof of \Cref{lem:oligo loc oligo}.}

\arka{Maybe we can send the proof of \Cref{lem:oligo loc oligo} to the appendix?}
%
\begin{proof}
If $(a_1,\dots,a_n)\in\orbit[\group]{(b_1,\dots,b_n)}$ then there exists $\gelem\in\group$ such that
\[
(\gelem\cdot a_1,\dots,\gelem\cdot a_n) =
\gelem\cdot(a_1,\dots,a_n) =
(b_1,\dots,b_n) \ .
\]
This means $\gelem\cdot a_i = b_i$ for all $i$.
This implies
\[
\gelem^{-1}\cdot b_i = \gelem^{-1}\cdot(\gelem\cdot a_i) =
(\gelem^{-1}\gelem)\cdot a_i = \iden[\group]\cdot a_i  = a_i\ .
\]
Hence $\gelem^{-1}\cdot(b_1,\dots,b_n) = (a_1,\dots,a_n)$,
which means $(b_1,\dots,b_n)\in\orbit[\group]{(a_1,\dots,a_n)}$
\end{proof}
%
\begin{proof}[Proof of \Cref{lem:oligo loc oligo}]
\begin{enumerate}
\item \Cref{lem:orbit part} implies that for every $n$ the set $\Indets^n$ partitions into disjoint orbits.
Since $\group\actson\Indets$ is oligomorphic we know this partition is finite.
Let $U_1,\dots,U_N$ be the orbits inside $\Indets^n$.
Pick $x_1,\dots,x_n\in\Indets$.
The set $V = \setof{(y_1,\dots,y_n)\in\Indets^n}{y_k\in\orbit[\monoid]{x_k}}$ is an equivariant subset of $\Indets^n$.
Therefore for every orbit $U_{\ell}$ inside $\Indets^n$ either $U_{\ell} \subseteq V$ or $U_{\ell} \cap V = \emptyset$.
Therefore $V$ is also a finite union of orbits. 
\item The only if direction follows from the previous point.
We prove the if direction.
Let $A\subseteqfin\Indets$ be a finite set such that $\orbit[\group]{A} = \Indets$.
Then $\Indets^n$ is the union of sets of the form $K(a_1,\dots,a_n)\defined\setof{(b_1,\dots,b_n)}{b_i\in\orbit[\group]{a_i}}$ where $(a_1,\dots,a_n)\in A^n$.
By \Cref{lem:orbit part} the set $\Indets^n$ partitions into orbits.
Each orbit inside $\Indets^n$ must be inside $K(a_1,\dots,a_n)$ for some $(a_1,\dots,a_n)\in A^n$.
There are finitely many such sets.
Hence to prove $\Indets^n$ is orbit-finite it is enough to show that each $K(a_1,\dots,a_n)$ is orbit-finite.
This follows from local oligomorphicity of $\group\actson\Indets$.
\end{enumerate}
\end{proof}
%
Now we give some examples and non-examples of locally oligomorphic monoid actions.
%
\begin{example}
The action $\inc{\OrderAtoms}\actson\OrderAtoms$ is effectively locally oligomorphic.
We show that it satisfies the three points in \Cref{def:loc oligo}.
%
\begin{enumerate}
\item For every $a\in\OrderAtoms$,
we have $\orbit[\inc{\OrderAtoms}]{a} = \OrderAtoms$.
Therefore, for any subset $S\subseteq\OrderAtoms$ of size $n$,
we have $\setof{(b_1,\dots,b_n)}{b_k\in\orbit[\monoid]{a_k}} = \OrderAtoms^n = \orbit[\inc{\OrderAtoms}]{S^n}$.
This proves $\inc{\OrderAtoms}\actson\OrderAtoms$ is locally oligomorphic.
\item For every $(a_1,\dots,a_n),(b_1,\dots,b_n)\in\OrderAtoms^n$ we have
$\orbit[\monoid]{(a_1,\dots,a_n)} = \orbit[\monoid]{(b_1,\dots,b_n)}$ if and only if $(a_1,\dots,a_n)$ and $(b_1,\dots,b_n)$ have the same order-type (for every $i,j$ we have $a_i \leq a_j$ if and only if $b_i \leq b_j$),
which is easy to check.
\item Since $\setof{(b_1,\dots,b_n)}{b_k\in\orbit[\monoid]{a_k}} = \OrderAtoms^n = \orbit[\inc{\OrderAtoms}]{S^n}$ for every $(a_1,\dots,a_n)\in\OrderAtoms^n$ and $S\subseteqfin\OrderAtoms$ of size $n$,
on any input we can simply output $\{1,\dots,n\}^n$.
\end{enumerate}
The output of the algorithm in the last point contains redundancies since two tuples in $\{1,\dots,n\}^n$ can be in the same orbit (for example, $(1,1,2)$ and $(2,2,3)$).
However, the redundancies can be removed simply by using the algorithm that allows us to check if two tuples are in the same orbit.
\end{example}
%
\begin{example}
Let $\calZ$ denote the structure of integers with the usual ordering and $\aut{\calZ}$ the group of all order preserving bijections of $\calZ$.
Then $\aut{\calZ}\actson\calZ$ is not locally oligomorphic.
Observe that $\orbit[\aut{\calZ}]{a} = \calZ$ for every $a\in\calZ$.
Therefore by \Cref{lem:oligo loc oligo} we have to just show that $\aut{\calZ}\actson\calZ$ is not oligomorphic.
For every $a,b,c,d\in\calZ$ the tuples $(a,b)$ and $(c,d)$ are in the same orbit if and only if $a - b = c - d$.
Therefore the set $\calZ^2$ splits into infinitely many orbits,
making $\aut{\calZ}\actson\calZ$ not oligomorphic.
\end{example}
%
\section{Equivariant Gröbner bases}
%
\arka{merge with next section}
%
\subsection{Monoid actions}
%
A \intro{monoid} $\monoid$ is a set equipped with a binary operation that is associative, and has an identity element.
In our setting, we are interested in polynomial rings of the form $\poly{\K}{\Indets}$ where $\Indets$ is an infinite set of indeterminates that are equipped with an injective \intro{monoid action} $\monoid \actson \Indets$.
This means that there is monoid homomorphism $\psi_{\monoid}$ from $\monoid$ to the monoid of one to one endofunctions on $\Indets$.
Obviously, the same monoid $\monoid$ can act differently on a set $\Indets$.
However, usually the action $\psi_{\monoid}$ will be clear from the choice of $\monoid$ and $\Indets$ and we will use $\monoid\actson\Indets$ to specify this choice.
For each $\gelem \in \monoid$, we use $\gelem\cdot x$ to denote $\psi_{\monoid}(\gelem)(x)$.
Usually, we will be interested in actions $\monoid\actson\Indets$ where $\Indets$ comes with some additional structure and $\monoid$ is the monoid of all structure preserving injective endofunctions/group of all structure preserving bijections on $\Indets$.
For instance, $\inc{\OrderAtoms}\actson\OrderAtoms$ denote the obvious action of the order preserving one-to-one endofunctions on the structure $\OrderAtoms$ of rational numbers with order.
In \arka{cite sec} we give some more examples of such actions.
%
\subsection{Equivariant ideals}
%
We use the notation $\poly{\K}{\Indets}$ for the ring of polynomials with coefficients from a field $\K$ and indeterminates from a set $\Indets$.
The action $\monoid\actson\Indets$ extends to an action $\monoid\actson\poly{\K}{\Indets}$ in the obvious way.
For illustration, $\pi\cdot (x^2 + 2y) = (\pi\cdot x)^2 + 2(\pi\cdot y)$.
A subset $S \subseteq \poly{\K}{\Indets}$ is called \intro{$\monoid$-equivariant} if for all $\gelem \in \monoid$ and $f \in S$,
we have $\gelem \cdot f\in S$.
An \intro{$\monoid$-ideal} of $\poly{\K}{\Indets}$ is an ideal of $\poly{\K}{\Indets}$ that is $\monoid$-equivariant.
We give in \Cref{ex:idl-equiv} examples and non-examples of \kl{$\monoid$-ideals}.
%
\begin{example}\label{ex:idl-equiv}
\arka{Notational convension for ideal used here}
The set $\idl[S]_0 \subset \poly{\K}{\Indets}$ of all polynomials whose set of coefficients sums to $0$ is an $\symgr{\Indets}$-ideal.
The set $\idl[M]_x$ of all polynomials that are multiple of a fixed $x \in \Indets$ is an \kl{ideal} that is not a $\symgr{\Indets}$-ideal.
\end{example}
%
\begin{proof}
We leave it to the reader to verify that $\idl[S]_0$ is indeed a $\symgr{\Indets}$-ideal.

To see that $\idl[M]_x$ is not a $\symgr{\Indets}$-ideal pick any variable $y\neq x$.
Let $\pi_{x,y}\in\symgr{\Indets}$ be the permutation that swaps $x$ and $y$ acts as identity everywhere else.
Then $\pi_{x,y}\cdot x = y \notin \idl[M]_x$.
\end{proof}

%
%\begin{proof}
%    Let $p,q\in \idl[S]_0$, and $r \in \poly{\K}{\Indets}$.
%    Then, $p \times r + q$ is in $S_0$. Remark that 
%    $p,r$ and $q$ belong to a subset $\poly{\K}{\Indets}$ of the 
%    polynomials that uses only finitely many indeterminates.
%    In this subset, the sum of all coefficients is obtained
%    by applying the polynomials to the value $1$ for every indeterminate
%    $y \in \Indets$. We conclude that
%    $(p \times r + q)(1,\dots, 1) 
%    = p(1,\dots,1) \times r(1,\dots,1) + q(1,\dots,1)
%    = 0 \times r(1, \dots, 1) + 0 = 0$, hence that
%    $p \times r + q$ belongs to $S_0$. 
%    Because $0$ is in $S_0$, we conclude that $S_0$ is an \kl{ideal}.
%    Furthermore, if $\gelem \in \group$ and $p \in S_0$, then
%    the sum of the coefficients $\gelem \cdot p$ is exactly
%    the sum of the coefficients of $p$, hence is $0$ too.
%    This shows that $S_0$ is \kl(ideal){equivariant}.
%
%    It is clear that all multiples of a given polynomial $x \in \Indets$
%    is an ideal of $\poly{\K}{\Indets}$. This is not an \kl{equivariant ideal}:
%    take any bijection $\gelem \in \group$ that does not map $x$ to $x$ (it
%    exists because $\Indets$ is infinite and $\group$ is all permutations),
%    then $\gelem \cdot x$ is not a multiple of $x$, and therefore does 
%    not belong to the ideal.
%\end{proof}
\AP
%
The orbit of an element $p$ in $\poly{\K}{\Indets}$ is defined as
\[
\intro*\orbit[\monoid]{p} \defined
\setof{\gelem\cdot p}{\gelem\in\monoid} \ .
\]
For $P\subseteq \poly{\K}{\Indets}$ we use $\intro*\orbit[\monoid]{P}$ to denote the union of the orbits $\orbit[\monoid]{p}$ for elements in $p$:
\[
\orbit[\monoid]{P} = \bigcup_{p\in P}\orbit[\monoid]{p} \ .
\]
An \kl{$\monoid$-equivariant set} is said to be \intro{orbit finite} if it is the union of finitely many \intro{orbits} under the action of $\monoid$,
equivalently, it is a set of the form $\orbit[\group]{P}$ for some finite set $P$.
Obviously, not every \kl{equivariant subset} is \kl{orbit finite}.
For instance, $\setof{x^n}{x\in\Indets, n\in\N} \subseteq \poly{\K}{\Indets}$ is $\symgr{\Indets}$-equivariant but is orbit-finite since $x^n$ and $y^m$ can not be in the same orbit when $n=m$.

The ideal \intro*{$\IdlGen{G}$} generated by a subset $G\subseteq \poly{\K}{\Indets}$ is the smallest ideal containing $G$.
For instance, the ideals $S_0$ and $M_x$ in \Cref{ex:idl-equiv} are generated by the sets $\setof{y - 1}{y\in\Indets}$ and $\{x\}$, respectively.
An ideal is called \intro{orbit-finitely generated} if it is generated by an orbit-finite set.
For example, the generating set $\setof{y - 1}{y\in\Indets}$ of $S_0$ forms a single orbit under $\symgr{\Indets}$.
Since orbit-finite sets are equivariant by definition,
orbit-finitely generated ideals are also equivariant.
Cohen showed in \cite{Cohen67} that every $\symgr{\Indets}$-equivariant ideal in $\poly{\K}{\Indets}$ is orbit-finitely generated.
However, as the next example shows,
this statement does not hold in general for every monoid actions.
%
\begin{example}
Consider the action $\monoid_0\actson\Indets$ where $\monoid_0$ is the trivial monoid,
i.e., a monoid with identity as its only element.
Since $\monoid_0$-orbits are singletons,
any ideal in $\poly{\Q}{\Indets}$ is an $\monoid_0$-ideal
and an ideal is orbit-finitely generated if and only if it is finitely generated.
However, if $\Indets$ is infinite,
there are ideals in $\poly{\Q}{\Indets}$ that are not finitely generated.

For instance, consider the ideal $\idl[I]$ of all polynomials in $\poly{\Q}{\Indets}$ without any constant term.
Any finite subset $F$ of $\idl[I]$ can not generate $\idl[I]$ since every non-zero polynomial in $\IdlGen{F}$ must use some variable that is used by some polynomial in $F$.
Therefore, $y\in\idl[I]\setminus\IdlGen{F}$ for any variable $y$ that is not used by any polynomial in $F$.
\end{example}
%
\begin{example}
Consider the action $\symgr{A}\actson A^2$ where $A$ is an countably infinite set defined as: $\pi\cdot (a,b) = (\pi\cdot a, \pi\cdot b)$.
We give an example of an $\symgr{A}$-equivariant ideal in $\poly{\Q}{A^2}$ which is not orbit-finitely generated.

For $n\geq 3$, define
\[
C_n \defined \setof{(a_0,a_1)(a_1,a_2)\dots(a_{n-2},a_{n-1})(a_{n-1},a_0) \in \poly{\Q}{A^2}}{a_i\neq a_j\in A}
\]
The set $C_n$ is an $\symgr{A}$-orbit for every $n$.
Let $\calC$ be the $\symgr{A}$-ideal generated by $\cup_{n \geq 3} C_n$.
Then $\calC$ does not have any orbit-finite set of generators.

\arka{give intuitive reason with graphs}

\arka{orbit $\mapsto$ $\monoid$-orbit}

\arka{$\monoid$-equivariant ideal $\mapsto$ $\monoid$-ideal?}

\end{example}
%
\begin{proof}
Pick an arbitrary orbit-finite subset $F$ of $\calC$.
We show $\IdlGen{F} \subsetneq \calC$.
Since the group $\symgr{A}$ acts in a variable-wise manner on $\poly{\Q}{A^2}$,
every polynomial in a $\symgr{A}$-orbit inside $\poly{\Q}{A^2}$ uses the same number of variables.
Since $F$ is orbit-finite,
there exists $N\in\N$ such that every polynomial in $F$ uses at most $N$-variables.
Therefore, $F\subseteq \IdlGen{\cup_{n = 3}^N C_n}$, and as a consequence $\IdlGen{F}\subseteq \IdlGen{\cup_{n = 3}^N C_n}$.
We show that for any $m > N$ we have $C_m \subseteq \calC \setminus \IdlGen{\cup_{n = 3}^N C_n} 
\subseteq \calC \setminus \IdlGen{F}$.

\arka{notation for monomials used here}

Pick $\monelt[p] \in C_m$ for some $m > N$.
If $\monelt[p]\in \IdlGen{\cup_{n = 3}^N C_n}$ then
\[
\monelt[p] = \sum_{i = 1}^k r_i \monelt[q]_i\monelt[p]_i
\]
for some $r_i\in\Q$, $\monelt[q]_i\in\mon{A^2}$ and $\monelt[p]_i\in \cup_{n = 3}^N C_n$.
The equality can hold only if $\monelt[p] = \monelt[q]_i\monelt[p]_i$ for some $i$.
This implies $\monelt[p]_i$ divides $\monelt[p]$.
We show that this leads to a contradiction.

Let
\[
\monelt[p] = (a_0,a_1)(a_1,a_2)\dots(a_{m-2}a_{m-1})(a_{m-1},a_0)
\]
and
\[
\monelt[p]_i = (b_0,b_1)(b_1,b_2)\dots(b_{n-2}b_{n-1})(b_{n-1},b_0) \ .
\]
For positive integers $k$, Let $\oplus_k$ denote the addition of integers modulo $k$.
Since $\monelt[p]_i$ divides $\monelt[p]$, $(b_0,b_1) = (a_j,a_{j\oplus_m 1})$ for some $j$.
But then $(b_1,b_2) = (a_{j\oplus_m 1},a_{j\oplus_m 2})$,
By induction, we must have $(b_{0\oplus_n\ell},b_{1\oplus_n\ell}) = (a_{j\oplus \ell},a_{(j\oplus_m (\ell+1)})$ for every $\ell$.
The LHS of this equality takes at most $m$-values but the RHS takes $n$.
Since $n > N \geq m$, this leads to a contradiction.
\end{proof}
%
\arka{define the notation $\EqIdlGen{H}$}.


\AP A function $f \colon X \to Y$ between two sets $X$ and $Y$ equipped with
actions $\group \actson X$ and $\group \actson Y$ is said to be
\intro(func){equivariant} if for all $\gelem \in \group$ and $x \in X$, we have
$f(\gelem \cdot x) = \gelem \cdot f(x)$. For instance, the
\kl(monomial){domain} of a monomial is an \kl{equivariant function}: if $\gelem
\in \group$, then $\gelem \cdot \dom(\monelt[m]) = \dom(\gelem \cdot
\monelt[m])$. Let us point out that the image of an \kl{orbit finite set} under
an \kl{equivariant function} is \kl{orbit finite}, and that the algorithms that
we will develop in this paper are all \kl(func){equivariant}.
%
\subsection{Equivariant Gröbner bases}
%
\subsubsection{Divisibility and reduction}

\AP
We use the notation $\poly{\K}{\Indets}$ for the ring of polynomials with coefficients from a field $\K$ and indeterminates from a set $\Indets$,
and $\mon{\Indets}$ for the set of monomials in $\poly{\K}{\Indets}$.
Letters $p,q,r$ denote polynomials,
$\monelt,\monelt[n]$ denote monomials, and $a,b,\alpha,\beta$ denote coefficients in $\K$.
\arka{reserve one of them for variables or check instances where $a,b$ are used for variables}


We know from \cite{GHOLAS24} that a
necessary condition for the \kl{equivariant Hilbert basis property} to hold is
that the set  $\mon{\X}$  of monomials is a \kl{well-quasi-ordering} when
endowed with the \intro{divisibility up-to $\group$} relation
($\intro*\gdivleq$), which is defined as follows: for $\monelt_1, \monelt_2 \in
\mon{\X}$, we write $\monelt_1 \gdivleq \monelt_2$ if there exists $\gelem \in
\group$ such that $\monelt_1$ \kl{divides} $\gelem \cdot \monelt_2$. This
relation also extends to monomials that are functions from $\Indets$ to
$(Y,\leq)$ with finite support, where $Y$ is any partially ordered set. We say
that a set $\Basis \subseteq \poly{\K}{\X}$ is an \intro{equivariant Gröbner
basis} of an equivariant ideal $\idl$ if $\Basis$ is \kl{equivariant},
$\IdlGen{\Basis} = \idl$, and for every polynomial $p \in \idl$, there exists
$q \in \Basis$ such that $\lm[\Indets](q) \gdivleq \lm[\Indets](p)$ and
$\dom(q) \subseteq \dom(p)$, following the definition of \cite{GHOLAS24}.

\AP We say that a group action $\group \actson \Indets$ is
\intro{well-structured} if the set $(\mon[Y]{\Indets}, \gdivleq)$ is
\kl{well-quasi-ordered}, for every \kl{well-quasi-order} $(Y,\leq)$. We say
that it is \intro{$\omega$-well-structured} if $(\mon{\Indets}, \gdivleq)$ is
\kl{well-quasi-ordered}.

Note that even in the case of a finite set of variables, a \kl{Gröbner basis}
is not necessarily an \kl{equivariant Gröbner basis}, because of the
\kl(polynomial){domain} condition. However, every \kl{equivariant Gröbner
basis} is a \kl{Gröbner basis}.

\begin{example}
  \label{ex:equivariant-gb}
  Let $\Indets \defined \set{ x_1, x_2 }$,
  with $x_1 \varleq x_2$,
  and $\group$ be the trivial group.
  Let us furthermore consider the ideal $\idl$ \kl(idl){generated by}
  $\set{ x_1, x_2 }$.
  Then, the set $\Basis \defined \set{ x_2 - x_1, x_1 }$ is a
  \kl{Gröbner basis} of $\idl$, but not an \kl{equivariant Gröbner basis}.
  Indeed, $x_2 \in \idl$, but there is no polynomial $q \in \Basis$
  such that $\lm(q) \divleq x_2$ and $\dom(q) \subseteq \dom(x_2)$.
\end{example}

\subsection{Computability of equivariant Gr\"{o}bner bases}
%
\arka{detailed description of our results}

We denote $\intro*\orbit[\monoid]{P}$ for the set of all elements $\gelem \cdot p$ for $\gelem \in \monoid$ and $p \in P$.
When $E$ is a singleton set ${p}$,
we write $\orbit[\monoid]{p}$ instead of $\orbit[\monoid]{\{p\}}$.
An \kl{equivariant set} is said to be \intro{orbit finite} if it is the
union of finitely many \intro{orbits} under the action of $\monoid$.
Equivalently, an \reintro{orbit finite set}
is a set of the form $\orbit[\group]{P}$ for some finite set $P$.
%
%
%
%\arka{Does not satisfy the domain condition}
%In the finite case, one can always compute an \kl{equivariant Gröbner basis} by
%computing \kl{Gröbner bases} for every possible ordering of the indeterminates,
%and taking their union.\footnote{This algorithm is correct because we are
%  considering the \kl{reverse lexicographic} ordering.}
 




